\section*{अध्याय \ref{ch:b3:genomics} --- जीनोमिकी और अनुक्रमण}

\begin{solution}{exo:b3:genomics:1}
किसी \omterm{def:b3:genomics:sanger}{डाइडिऑक्सीन्यूक्लियोटाइड} में 3$'$ हाइड्रॉक्सिल नहीं होता, इसलिए
अगले न्यूक्लियोटाइड से कोई फ़ॉस्फ़ोडाइएस्टर बंध नहीं बन सकता और शृंखला
रुक जाती है। केवल डाइडिऑक्सी रूप हों तो हर शृंखला पहली ही स्थिति पर रुक
जाएगी; केवल सामान्य रूप हों तो कोई नहीं रुकेगी। मिश्रण समापन को हर
स्थिति पर एक यादृच्छिक घटना बना देता है, इसलिए उत्पाद एक पूरी सीढ़ी
बनाते हैं, साँचे के हर क्षार के लिए एक लंबाई।
\end{solution}

\begin{solution}{exo:b3:genomics:2}
$4\times 10^{8}\times \qty{300}{bp} = 1.2\times 10^{11}$ क्षार। मनुष्य:
$1.2\times 10^{11}/3.2\times 10^{9} \approx 38\times$। जीवाणु:
$1.2\times 10^{11}/5\times 10^{6} = \num{24000}\times$।
\end{solution}

\begin{solution}{exo:b3:genomics:3}
कॉन्टिग अतिव्यापी \omterm{def:b3:genomics:ngs}{वाचनों} से संग्रथित कोई सतत अनुक्रम है; \omterm{def:b3:genomics:assembly}{स्कैफ़ोल्ड}
कॉन्टिगों का ऐसा क्रमित, दिशित समुच्चय है जिनके बीच आकलित आकार के
अंतराल हैं; \omterm{def:b3:genomics:assembly}{N50} वह कॉन्टिग लंबाई है जिस पर छाँटे गए कॉन्टिग संग्रथित
क्षारों के आधे तक पहुँच जाते हैं। यहाँ अकेले \qty{8}{Mb} के दस कॉन्टिग
में \qty{80}{Mb} हैं, जो \qty{50}{Mb} के आधे बिंदु से आगे है (सातवाँ
\qty{56}{Mb} तक पहुँच जाता है): \omterm{def:b3:genomics:assembly}{N50} $= \qty{8}{Mb}$।
\end{solution}

\begin{solution}{exo:b3:genomics:4}
जीनोम का आकार जीव की जटिलता के साथ नहीं चलता: किसी प्याज़ में मनुष्य से
पाँच गुना DNA है, किसी फेफड़ा-मछली में चालीस गुना; खमीर और कृमि के DNA
में दस गुना अंतर है पर जीन-संख्याएँ मिलती-जुलती हैं। अतिरिक्त भाग
अकूटलेखी है: स्थानांतरणीय अवयव और उनके अवशेष, अंतर्वेशन, उपग्रह
पुनरावृत्तियाँ।
\end{solution}

\begin{solution}{exo:b3:genomics:5}
$e^{-c} = 10^{-6}$ से $c = 6\ln 10 = 13.8$ मिलता है। $c = 8$ के लिए: $N =
cG/L = 8\times 5\times 10^{6}/150 \approx \num{267000}$ \omterm{def:b3:genomics:ngs}{वाचन}, $\theta
= 30/150 = 0.2$, कॉन्टिग $= N e^{-6.4} = \num{267000}\times 0.00166
\approx 440$।
\end{solution}

\begin{solution}{exo:b3:genomics:6}
$P(X < 8) = P(X \le 7) \approx P\bigl(Z < (7.5 - 15)/2.74\bigr) =
P(Z < -2.74) \approx 0.003$। $30\times$ पर किसी विषमयुग्मजी के दोनों
युग्मविकल्पी लगभग सदा कई बार दिखते हैं, आच्छादन असमान होता है (GC-समृद्ध
क्षेत्रों को कम \omterm{def:b3:genomics:ngs}{वाचन} मिलते हैं, इसलिए कुछ स्थलों पर औसत का आधा ही
दिखता है), और इतने \omterm{def:b3:genomics:ngs}{वाचन} बचे रहते हैं कि किसी असली युग्मविकल्पी को
$10^{-3}$ त्रुटियों से अलग किया जा सके।
\end{solution}

\begin{solution}{exo:b3:genomics:7}
पुनरावृत्ति के भीतर से आया हर \qty{150}{bp} \omterm{def:b3:genomics:ngs}{वाचन} एक जैसा होता है, चाहे
वह किसी भी प्रति से आया हो, इसलिए संग्रथक एक ही \qty{6}{kb} शीर्ष
बनाता है जिसमें $500$ पथ घुसते और $500$ निकलते हैं; वह बता नहीं सकता कि
कौन-सा प्रवेश किस निकास से जुड़ता है, और \omterm{def:b3:genomics:assembly}{संग्रथन} $500$ अंतरालों में टूट
जाता है, जिसमें पुनरावृत्ति एक ही बार, माध्य से $500$ गुना आच्छादन के
साथ, उपस्थित रहती है। पुनरावृत्ति और दोनों ओर के अनन्य पार्श्वों से
लंबे \omterm{def:b3:genomics:ngs}{वाचन} --- \qty{8}{kb} या अधिक --- हर प्रति को पाट देते हैं और उसे
सुलझा देते हैं।
\end{solution}

\begin{solution}{exo:b3:genomics:8}
$4.8\times 10^{7}/1000 = \num{48000}$ कूटलेखी अंतर; एक-तिहाई, लगभग
\num{16000}, प्रोटीन बदलते हैं।
\end{solution}

\begin{solution}{exo:b3:genomics:9}
प्रत्याशित मिथ्या धनात्मक $10^{6}\times 5\times 10^{-8} = 0.05$।
\qty{2}{\%} वाले रोग पर $1.15$ का सापेक्ष जोखिम प्रति हज़ार कुछ ही
स्थितियाँ बदलता है; हज़ार लोगों में लगभग बीस रोगी होते हैं, जो उसे
देखने के लिए कहीं कम हैं (सामर्थ्य रोगियों की संख्या और प्रभाव के वर्ग
के साथ बढ़ती है)। किसी व्यक्ति के लिए $0.3$ प्रतिशत-बिंदु निरर्थक है।
पर वह युग्मविकल्पी ऐसे जीन को चिह्नित करता है जिसकी क्रियाशीलता में
मामूली बदलाव रोग-जोखिम बदल देता है --- वह जीन, और उसका पथ, ऐसा
औषध-लक्ष्य हो सकता है जिसके पूर्ण निरोध का बड़ा प्रभाव हो।
\end{solution}

\begin{solution}{exo:b3:genomics:10}
समष्टि-आनुवंशिक: विशाल समष्टियों वाले जीवाणुओं में चयन थोड़े महँगे
अंतर्वेशन भी हटा देता है; छोटी प्रभावी समष्टियों वाले स्तनधारियों में
अपवाह हल्के हानिकारक अकूटलेखी DNA को जमा होने देता है --- वंशक्रमों के
पार जीनोम-आकार और समष्टि-आकार के बीच विलोम सहसंबंध से समर्थित, अपवादों
से कमज़ोर। संरचनात्मक: सुकेंद्रकी जीनोमों पर वे ट्रांसपोज़ॉन आक्रमण
करते हैं जिन्हें जीवाणु, विलोपन की अपनी अभिनति और स्वार्थी अवयवों के
लिए किसी अर्धसूत्री शरण के अभाव में, साफ़ कर देते हैं --- जीनोम-आकार और
ट्रांसपोज़ॉन-अंतर्वस्तु के सहसंबंध से समर्थित। नियामक: जटिल परिवर्धन को
अधिक नियामक DNA चाहिए --- परिवर्धन-जीनों के चारों ओर संरक्षित अकूटलेखी
अवयवों से समर्थित, पर इस तथ्य से कमज़ोर कि मानव जीनोम का केवल
\qty{8}{\%} ही कोई चयन दिखाता है, इसलिए अधिकांश अकूटलेखी DNA नियामक
नहीं है।
\end{solution}

\begin{solution}{exo:b3:genomics:11}
दोनों प्रकार के आरंभ कुल घनत्व $\rho = (N_{1} +
N_{2})/G$ से पड़ते हैं। लंबाई $L_{i}$ का कोई \omterm{def:b3:genomics:ngs}{वाचन} अपने कॉन्टिग का सबसे
दायाँ तब है जब उसके बाद के $L_{i} - T$ क्षारों में किसी भी प्रकार का
कोई \omterm{def:b3:genomics:ngs}{वाचन} आरंभ न हो: प्रायिकता $e^{-\rho(L_{i} - T)}$। इसलिए कॉन्टिग $= N_{1}
e^{-\rho(L_{1} - T)} + N_{2} e^{-\rho(L_{2} - T)}$। कोई दीर्घ \omterm{def:b3:genomics:ngs}{वाचन}
किसी छोटे \omterm{def:b3:genomics:ngs}{वाचन} की तुलना में चरघातांकी रूप से कम प्रायिकता से सबसे दायाँ
होता है, और हर दीर्घ \omterm{def:b3:genomics:ngs}{वाचन} अपनी पूरी लंबाई पर किसी अंतराल की संभावना भी
हटा देता है; छोटे \omterm{def:b3:genomics:ngs}{वाचनों} के रूप में उतने ही क्षार अनेक आरंभ-बिंदु
जोड़ते हैं पर हर सुरक्षित खिड़की छोटी होती है, इसलिए दीर्घ \omterm{def:b3:genomics:ngs}{वाचन} जीतते हैं।
\end{solution}

\begin{solution}{exo:b3:genomics:12}
दो \omterm{def:b3:genomics:comparative}{पूर्ण-जीनोम द्विगुणन} यह पूर्वानुमान करते हैं कि चौगुना प्रतिरूप
जीनोम भर मिलेगा: पैरालॉगी गुणसूत्र-खंडों के चतुष्क (पैरालॉगॉन) जो एक ही
क्रम में एक ही जीन-कुल ढोते हों, और जिनके द्विगुणन सारे कुलों के लिए
एक ही काल के हों; ऐम्फ़िऑक्सस और \emph{Ciona} में एक ही गुच्छा, जो
घटनाओं से पहले अलग हुए; और लैम्प्रे में, जो उनके आसपास अलग हुए, कोई
मध्यवर्ती या स्वतंत्र रूप से बनी अवस्था। अकेले Hox गुच्छे के स्वतंत्र
द्विगुणन केवल Hox के लिए पैरालॉगॉन, भिन्न इतिहास वाले पड़ोसी, और कुलों
के बीच भिन्न द्विगुणन-काल पूर्वानुमानित करते हैं। जीनोम-व्यापी
पैरालॉगॉन और साझा काल-निर्धारण, जैसा प्रेक्षित है, \omterm{def:b3:genomics:comparative}{पूर्ण-जीनोम
द्विगुणन} के पक्ष में हैं।
\end{solution}

\begin{solution}{pb:b3:genomics:1}
\textbf{1.} $0.002\times 6\times 10^{8} = 1.2\times 10^{6}$ \omterm{def:b3:genomics:ngs}{वाचन};
$c = 1.2\times 10^{6}\times 150/4.6\times 10^{6} \approx 39$।
\textbf{2.} $e^{-39} \approx 10^{-17}$: प्रभावी रूप से कोई क्षार
अनाच्छादित नहीं।
\textbf{3.} $\theta = 0.2$; कॉन्टिग $= 1.2\times 10^{6} e^{-31} \approx
0$: सिद्धांततः एक कॉन्टिग, अंतराल केवल पुनरावृत्तियों पर।
\textbf{4.} $c = \num{30000}\times 150/4.6\times 10^{6} = 0.98$;
अनाच्छादित $e^{-0.98} = 0.38$, लगभग \qty{1.7}{Mb}; कॉन्टिग $= \num{30000}
\times e^{-0.78} \approx \num{13700}$।
\textbf{5.} सातों प्रकार्यगुच्छ एक जैसे \omterm{def:b3:genomics:ngs}{वाचन} देते हैं और एक ही
\qty{5}{kb} कॉन्टिग में ढह जाते हैं, जिसमें सात अनन्य बाएँ पार्श्व
घुसते हैं और सात दाएँ पार्श्व निकलते हैं और उनका युग्मन अज्ञात रहता है:
$14$ कॉन्टिग-सिरे, सात अंतराल।
\textbf{6.} लगभग \num{4600} जीन; कूटलेखी $0.88\times 4.6\times 10^{6}
= 4.0\times 10^{6}$ bp, माध्य जीन \qty{880}{bp}।
\textbf{7.} $0.998\times 6\times 10^{8}\times 150/3.2\times 10^{9}
\approx 28\times$।
\textbf{8.} प्रति युग्मविकल्पी लगभग $14$ \omterm{def:b3:genomics:ngs}{वाचन}।
\textbf{9.} किसी दिए हुए गलत क्षार को दिखाते $28\times 10^{-3}/3 \approx 0.009$
\omterm{def:b3:genomics:ngs}{वाचन} --- तीन या अधिक की संभावना लगभग $10^{-7}$ है --- और
$28$ का \qty{20}{\%} $5.6$ \omterm{def:b3:genomics:ngs}{वाचन} है; कोई त्रुटि दोनों में से कोई
परीक्षण पार नहीं करती, जबकि कोई असली विषमयुग्मजी युग्मविकल्पी $14$ में
प्रत्याशित है।
\textbf{10.} $3.2\times 10^{9}/1500 \approx 2.1\times 10^{6}$
विषमयुग्मजी स्थल; कूटलेखी अनुक्रम में $4.8\times 10^{7}/1500 \approx
\num{32000}$।
\textbf{11.} किसी Alu से आया \qty{150}{bp} \omterm{def:b3:genomics:ngs}{वाचन} $1.1$ करोड़ प्रतियों
में से अनेक से लगभग बराबर अच्छा मेल खाता है, इसलिए वह गलत प्रति पर रख
दिया जाता है या उसे कम संरेखण-विश्वास मिलता है। तब प्रतियों के बीच के
अंतर विषमयुग्मजी प्रकारांतरों जैसे दिखते हैं, और असली प्रकारांतर उनके
बीच छिप जाते हैं: Alu के भीतर के प्रकारांतर-निर्धारण प्रायः छान दिए
जाते हैं, और ये अवयव छोटे-वाचन अनुक्रमण के अंधे क्षेत्र हैं।
\textbf{12.} $1.2\times 10^{-8}\times 6.4\times 10^{9} \approx 77$ नए
उत्परिवर्तन। बच्चे में लगभग $14$ \omterm{def:b3:genomics:ngs}{वाचनों} को प्रकारांतर दिखाना चाहिए; पर
यदि किसी जनक का स्थल केवल पाँच \omterm{def:b3:genomics:ngs}{वाचनों} से ढका हो तो कोई विषमयुग्मजी
युग्मविकल्पी प्रायिकता $2^{-5} = 3\,\%$ से छूट जाता है और कोई वंशागत
प्रकारांतर नया मान लिया जाता है --- इसलिए जनकों का अनुक्रमण बच्चे जितना
ही गहरा होना चाहिए।
\textbf{13.} $4.8\times 10^{7}\times 100/150 = 3.2\times 10^{7}$ \omterm{def:b3:genomics:ngs}{वाचन},
जीनोम से बीस गुना कम: अनुक्रमण की लागत बीस गुना गिरती है, जो पकड़ने के
चरण की लागत से अधिक है।
\textbf{14.} $1.1\times 10^{6}\times 300 = 3.3\times 10^{8}$ bp, यानी
जीनोम का लगभग \qty{10}{\%}; \omterm{def:b3:genomics:ngs}{वाचनों} का \qty{10}{\%}, कोई $6\times
10^{7}$।
\textbf{15.} $\num{60000}\times 1500 = 9\times 10^{7}$ bp; $9\times
10^{7}/1.6\times 10^{10} = 0.56\,\%$।
\textbf{16.} $0.012\times 2.9\times 10^{9} = 3.5\times 10^{7}$
प्रतिस्थापन, प्रति वंशक्रम $1.7\times 10^{7}$; दर $1.7\times
10^{7}/(2.9\times 10^{9}\times 6.5\times 10^{6}) = 9\times 10^{-10}$
प्रति क्षार प्रति वर्ष, 25 वर्ष की प्रति पीढ़ी $2.3\times 10^{-8}$ ---
वंशावली-दर $1.2\times 10^{-8}$ से लगभग दोगुना, एक ज्ञात विसंगति जो
अतीत में लंबे पीढ़ी-काल या अधिक पुराने विभाजन की ओर इशारा करती है।
\textbf{17.} \num{16000} की चार प्रतियाँ \num{64000} होतीं;
\num{20000} बचे हैं, इसलिए \num{48000} अतिरिक्त प्रतियों में से केवल
\num{4000} बची हैं: द्विगुणों का \qty{92}{\%} खो गया।
\textbf{18.} कृमि और मनुष्य में जीनों की संख्या समान है; जटिलता इसमें
है कि उनका उपयोग कैसे होता है --- वैकल्पिक विच्छेदन एक जीन को हज़ारों
प्रोटीनों में गुणा कर देता है (\emph{Dscam} के लिए \num{38016}), नियमन
अनुलेखन-कारकों और संवर्धकों को कोशिका-प्रकार-विशिष्ट ढंग से जोड़ता है,
\omterm{def:b3:rna-regulation:ncrna}{अकूटलेखी RNA} परतें जोड़ते हैं, और प्रोटीन जालों में अन्योन्यक्रिया करते हैं।
\textbf{19.} नियामक अवयव (संवर्धक, प्रवर्तक, विद्युतरोधक), \omterm{def:b3:rna-regulation:ncrna}{अकूटलेखी RNA}
जीन, विच्छेदन तथा स्थानिकीकरण संकेत, उद्गम और संरचनात्मक अनुक्रम। किसी
संभावित संवर्धक का परीक्षण उसे किसी पारजीनी भ्रूण में किसी प्रतिवेदक
जीन के आगे रखकर और यह देखकर कीजिए कि प्रतिवेदक कहाँ अभिव्यक्त होता है,
फिर उस अवयव को जीनोम में CRISPR से विलोपित करके पड़ोसी जीनों को मापिए।
\textbf{20.} $0.05/10^{6} = 5\times 10^{-8}$; $p < 10^{-5}$ पर $10^{6}
\times 10^{-5} = 10$ मिथ्या धनात्मक प्रत्याशित।
\textbf{21.} अवाहक संभावना $0.009/0.991 = 0.00908$; एक प्रति संभावना
$0.00954$, जोखिम $0.945\,\%$; दो प्रतियाँ संभावना $0.01001$, जोखिम
$0.99\,\%$।
\textbf{22.} अंक $200$ युग्मविकल्पियों पर यह जोड़ता है कि कोई व्यक्ति
कितनी जोखिम-प्रतियाँ ढोता है, हर युग्मविकल्पी के लघुगणकीय
संभावना-अनुपात से भारित। गणना का माध्य $200\times 2\times 0.3 = 120$ है
और मानक विचलन $\sqrt{200\times 2\times 0.3\times 0.7} \approx 9$;
ऊपरी \qty{1}{\%} औसत से लगभग $21$ युग्मविकल्पी अधिक ढोता है, और
$1.05^{21} \approx
2.8$: औसत व्यक्ति की संभावना से लगभग तीन गुना, उन युग्मविकल्पियों से जो
अकेले जोखिम \qty{5}{\%} बदलते हैं।
\textbf{23.} दसियों किलोक्षार के किसी खंड के भीतर के युग्मविकल्पी साथ
वंशागत होते हैं (सहलग्नता असंतुलन), इसलिए उनमें से कोई भी वही सहचर्य
दिखाता है; जिस SNP का जीन-प्ररूपण हुआ वह बस वही है जो पट्टिका पर था।
कारक प्रकारांतर ढूँढ़ने के लिए: रोगियों और नियंत्रणों में उस क्षेत्र का
अनुक्रमण कीजिए, समुच्चय को सांख्यिकीय रूप से सँकरा कीजिए, फिर हर
संभावित प्रकारांतर का परीक्षण कीजिए --- हर युग्मविकल्पी की संवर्धक
क्रियाशीलता के लिए प्रतिवेदक परीक्षण, कोशिकाओं में हर युग्मविकल्पी का
संपादन करके पास के जीनों की अभिव्यक्ति मापना, अभिव्यक्ति-मात्रात्मक-लक्षण
संकेतों के साथ सह-स्थानिकीकरण।
\textbf{24.} जीनोम का $1 - e^{-0.5} = 39\,\%$ पढ़ा जाता है। वह अतीत के
किसी ज्ञात काल पर किसी समष्टि का सीधा प्रतिदर्श देता है --- बाद के
प्रवसनों और मिश्रण से पहले की युग्मविकल्पी बारंबारताएँ, नियंडरथाल खंडों
की लंबाई (लंबी, क्योंकि पुनर्संयोजन की थोड़ी ही पीढ़ियों ने उन्हें तोड़ा
था) और इससे मिश्रण की तिथि --- जिसका अनुमान आधुनिक जीनोम केवल
प्रतिरूपों से लगा सकते हैं।
\textbf{25.} प्रायोगिक चरण में लगभग \num{13700} कॉन्टिग; प्रति अगुणित
जीनोम $28\times$, प्रति युग्मविकल्पी लगभग $14$ \omterm{def:b3:genomics:ngs}{वाचन}; किसी बच्चे में
कोई $77$ नए उत्परिवर्तन।
\end{solution}
